\documentclass{article}
\usepackage{amsmath,amssymb,hyperref}
\title{Wait for the named permission and the consuming ledger, not a forecast of both}
\date{26 November 2026}
\begin{document}
\maketitle
\paragraph{Finding.}
An earlier transmission need not keep a replacement report out until
a predicted absolute time. A local decision can instead require both
the qualified physical handoff receipt and enough room on the
\emph{actual consuming ledgers}. This avoids needing the remote motion
clock to compute budget age, but requires the authorized ability to
hold at the local broker and an upper bound on the time usable evidence
arrives. It may add delay to the old report. An actual retry without
capacity can consume sender budget even if its recipient rejects it;
defer before charging, not after. These are source-fixture outcomes,
not claims of live use, policy superiority over a matched direct
scheduler or I1/I2 discharge.

\paragraph{Which clock and which physical event.}
Use byte-pinned native warehouse steps as in
\href{https://github.com/eternal-november/eos-research-programmes/pull/625}{the
departure-qualified binding replay}. The simulated appointment occurs only
after the named carrier is unloaded and reaches the named exit.
Native motions place this at ticks \(4,5,6\) in the finite envelope.
The experiment feeds only an arrival cut \(r\), in the consuming
broker's eighth-second grid, to the replacement handler.
An authored bridge signs scope in the model and sometimes delays or
omits qualified receipt. No stock warehouse or broker is claimed to
implement that bridge. It is not legal to act on a remote event merely
because the evaluator knows it occurred. Within the fixture, once
received, the exact permission and chosen physical receiver remain valid
through the selected dispatch or deadline. If the law allows revocation
or expiry during the hold, the owner must check that law too; the one
lower threshold below would no longer suffice.

The native SubT radio uses the same actual sender and receiver debit
objects, started from clean ledgers in the fixture. It debits the
sender on an admitted transmission even if the optional peer's route
erases it, while the receiver is only debited on successful arrival.
Actual application callback is the measure of paid here.
Either shared sender or shared physical receiver can limit the old
report. Logical topic identity alone is no proxy for physical ledger.
All operators are compliant, and this is a controlled scheduling
experiment with integer byte budgets, not a network security test.

\paragraph{Exact one-predecessor rule.}
Let \(B\) have been transmitted at local cut zero. Let \(A\) be a
named live duty, whose single individually fitting send has a valid
local receipt at \(r\) and is owed by the boundary \(D\).
Use the admitted reliable path for \(A\) and no additional blocker
between receipt and use. Write \(P\) when some actual shared sender
or receiver carries a still-relevant debit from B that with A's
demand would exceed that row's cap. Under permanent unchanged scope,
\(P\) computed from the early debit is an initial conflict, not a
fresh status guessed for later. A direct check at actual consumption
counts the active rows and so need not infer \(P\) from the previous
reply. The discrete model admits \(A\) precisely at
\[
 t\ge r \quad\text{and}\quad (\neg P\;\text{or}\;t>8).
\]
Hence first feasible dispatch is \(r\) without conflict, or
\(\max(r,9)\) with conflict. The old row remains active at \(t=8\)
exactly. These cuts count one eighth of a second each.
The result follows since a lone prior debit is either live in
its own row or fully expired, with no debit ever duplicated.
For a finite predeclared future support before committing B,
the required positive guarantee is that the earliest cut
is at most \(D\) in every admitted history. If an admissible
receipt is lost or late beyond \(D\), correct refusal alone
does not fulfill the duty. When the worst local receipt bound and
the relevant local expiry bound are both within \(D\), the
guard needs no exact translation of remote event timestamps.
This still assumes a justified bound on qualified evidence
availability and a real dispatch opportunity at these cuts.
It does not retime an old ledger or a cross-clock deadline.

A conservative arm waits out the time-zero transmission in
all cases, with the same permission and resources. Even if B
was erased and could affect only the recipient row, it will
wait through cut 8. The exact arm can serve A at cut 8 in that
case, since B spent no shared receiver debit, while it still
waits if they share the sender. Independent same-information
native-ledger enumeration reaches the same rule and matches.

\paragraph{Actual consumption and another live duty.}
A separate experiment tests the need to check at the actual
consuming cut. At native cut 10 (1.25 seconds), a separately
authorized mandatory arrival \(C\) is processed on receiver R1.
A prior check for A had found room before C. Then D, an
independently owed report on the same sender as A but a
different actual receiver, must be dispatched at cut 10 too.
Compare:
\begin{enumerate}
\item A stale check proceeds as if approval reserved capacity:
C runs, A attempts using the old answer, then D tries.
\item An authorized direct receiver commits the mandatory C and
D cuts, and rechecks A at the same actual consuming ledger.
It allows the first timely fitting slot and otherwise refuses
without attempting A or charging its sender.
\end{enumerate}
Both C and D have physically named native application outputs.
For common cap five and each payload size three, if A uses
R1 the stale A actually fails at receive, after debiting three
TX bytes. D then fails too even though its receiver is disjoint.
A fresh serialized check preserves the due D callback. It
also preserves A at cut 19 if its latest due cut is 19; at
18 it correctly refuses and reports the old duty unfulfilled.
With fixed C before the other calls and D due at the same
cut, reversing the A/D order cannot make both fit if their
sender sum exceeds cap. If all row sums fit, either order
delivers. Joint service after the cut therefore needs the
known release time or another independently lawful capability.
Protecting A by forbidding all C/D work is not an admissible
resource reservation in this scenario. A scope-matched
priority/serialization implementation with identical information
and the same right to defer is an ordinary strong comparator
and ties, not a weaker algorithm.

The check-plus-attempt is serialized with consuming writers
in the selected harness. An unsynchronized read plus a
later command is not an interlock, and the study does not
establish that a stock online API exposes such a consuming
boundary. Other sends, clock stalls, endpoint edits or failures
during the hold require a new envelope. There is no second
attempt laundered as completed old service in the stale arm.
The local exact arm never sends A when its relevant current
row would already reject it. It does not turn the no-room
status into confirmation that a prior effect settled.

\paragraph{Audit.}
Across 12,960 native cases per one-debit arm, the exact
consumer meets 6,285 old deadlines and the unconditional-age
hold 5,832. A case without a still-valid receipt is not
counted as a delivery. The 4,320 support checks compare
preflight deadlines with direct enumeration. In 486 same-cut
ordering cases per arm, the stale check incurs 117 failed old
attempts and 252 actual side deliveries. The serialized arm
has no failed old attempts and 486 side deliveries; it
completes the old duty in 396 cases. The direct lawful
same-resource search agrees on the earliest feasible cuts.
These are controlled case counts, not estimated field
reliability. Lean checks the exact deadline/existence and
row-overflow lemmas. Receiver evidence from real communication,
enforcement through actual effect under further failures,
strongest independent mission comparisons, and contract
acceptance are separate and remain open.
\end{document}
